\section{Related Work}
\label{sec:related_work}

\subsection{Residual feedback for multi-horizon forecasts}
Prediction errors from completed forecasts can inform the next correction. ResCAL~\cite{rescal2022} estimates future traffic-forecasting errors from previously observed residuals and graph signals. TEFL~\cite{tefl2026} addresses the timing of residual feedback in rolling multi-horizon forecasts: it selects the most recent fully observable residual vector, maps that vector through a low-rank additive adapter, and jointly trains the adapter and base forecaster after a base-model warmup. TEFL's complete residual input already includes its final value. We examine how to retain that value explicitly when the residual block is compressed.

EPOC addresses the retained representation of a completed block. It keeps a small set of fixed DCT coefficients together with the uncompressed final residual, shares that endpoint across component-wise regressions, and also uses coefficients of the current base forecast. The base forecaster remains fixed while the regression statistics update after each completed block. The joint-base comparison in Table~\ref{tab:learned_comparison} evaluates this correction against a TEFL-style residual adapter; the equal-size summary in Table~\ref{tab:endpoint_selection} and input controls in Fig.~\ref{fig:input_ablation} examine the retained endpoint and its combination with compressed features.

\subsection{Calibration of fixed forecasters}
Methods that retain a fixed base differ in the information supplied to their auxiliary predictors and in what they update. ELF~\cite{elf2025} learns an online forecast from Fourier features of observed history and combines it with the fixed model's forecast through an adaptive weight. TAFAS~\cite{tafas2025} places gated temporal calibration modules around the base and can adapt using partially observed targets. PETSA~\cite{petsa2025} uses low-rank input and output calibration modules with robust, spectral, and patch-structure loss terms. The $\delta$-Adapter~\cite{deltaadapter2026} attaches bounded input and output adjustments to a deployed forecaster, including an additive output correction.

COSA~\cite{cosa2026} applies a gated linear correction to the current forecast using statistics of recently observed targets as context. FAC~\cite{fac2026} calibrates the input and output through gated complex-affine transformations in the frequency domain, updating from matured targets. OMPB~\cite{ompb2026} trains a gated Bayesian residual head on a fixed forecast under an online objective that includes source-risk and source-to-target mismatch terms. These methods establish several routes to online correction: a separate forecast, input/output calibration, or an output-side adapter. Our correction uses the preceding forecast's residual coefficients and endpoint as explicit inputs, with the current forecast coefficients, and constrains its output to a fixed DCT subspace. Table~\ref{tab:prior_methods} compares the auxiliary information, feedback mechanism, and retained state of the methods evaluated in this study.

\subsection{Sequential adaptation and retained information}
Other online forecasters adapt the forecasting system or regulate when it should change. FSNet~\cite{fsnet2023} combines fast adaptation with memory for recurring patterns, and OneNet~\cite{onenet2023} updates and combines models of temporal and cross-variable dependencies. Reconditionor identifies context-driven shifts through residual--context dependence, while SOLID selects similar samples to fine-tune a prediction layer~\cite{solid2024}. DynaTTA~\cite{dynatta2025} adjusts adaptation using prediction-error and representation-drift signals; PADRE~\cite{padre2026} uses delayed residual trajectories and retrieved drift patterns to gate backbone updates. EPOC keeps the backbone fixed and uses completed residuals directly as features for an auxiliary correction.

The DCT~\cite{ahmed1974dct}, ridge regression~\cite{hoerl1970ridge}, and forgetting in sequential least squares~\cite{goel2020rls} provide established tools for compressing and estimating from a stream. The design choice examined here is which information from a completed residual block to keep alongside its compressed coefficients. The endpoint, residual coefficients, and current-forecast coefficients define the regression input; the fixed output subspace and sufficient statistics determine the retained state analyzed in Section~\ref{sec:method}.

\begin{table*}[t]
\caption{Auxiliary inputs, correction mechanisms, and retained state in the evaluated comparison methods. Existing methods are ordered by publication year; EPOC appears last.}
\label{tab:prior_methods}
\centering
% Focused structural comparison; numerical transfers are specified in Section IV.
\begingroup\fontsize{9}{11}\selectfont
\setlength{\tabcolsep}{3pt}
\begin{tabular}{@{}>{\raggedright\arraybackslash}p{.20\linewidth}>{\raggedright\arraybackslash}p{.24\linewidth}>{\raggedright\arraybackslash}p{.28\linewidth}>{\raggedright\arraybackslash}p{.20\linewidth}@{}}
\toprule
Method & Auxiliary input & Correction and feedback & Retained state \\
\midrule
ELF (ICML 2025) & Fourier features of observed history & Online auxiliary forecast and adaptive mixing & Regression and mixing statistics \\[3pt]
TEFL (arXiv 2026) & Most recent completed residual block & Low-rank residual adapter; joint base training & Adapter weights and residual input \\[3pt]
$\delta$-Adapter (ICLR 2026) & Current input and forecast & Bounded input/output adapters; online updates & Adapter weights and optimizer state \\[3pt]
COSA (ICLR 2026) & Current forecast and recent target statistics & Gated linear output correction; online updates & Linear adapter, gate, context \\[3pt]
FAC (arXiv 2026) & Frequency representations of input and forecast & Gated frequency-domain calibration & Masks, biases, optimizer state \\[3pt]
OMPB (UAI 2026) & Base forecast & Gated Bayesian residual head; predict then update & Head, optimizer, source/target buffers \\[3pt]
EPOC & Past residual DCT coefficients and endpoint; current forecast coefficients & Fixed DCT output subspace; complete-block update & Ridge statistics, endpoint, mixing state \\
\bottomrule
\end{tabular}
\endgroup

\end{table*}
